\section{Evaluation Framework}
\label{sec:4_8}

Our evaluation framework is designed to provide a comprehensive assessment of model performance through various metrics, temporal validation, and baseline comparisons. It also incorporates real-time monitoring through drift detection algorithms. The implementation details can be found in the \texttt{src/evaluation/} directory and the \texttt{src/utils/metrics.py} file.

\subsection{Performance Metrics}
\label{sec:4_8_1}

The framework calculates several metrics to capture different aspects of how the model performs:

\paragraph{Primary Metric: AUC-PR}
We use the Area Under the Precision-Recall Curve (AUC-PR) as our primary metric because of the severe class imbalance in our data (an 8.11% fraud rate). Unlike the more common AUC-ROC, which can be overly optimistic when one class dominates, AUC-PR directly measures the trade-off between precision and recall across all classification thresholds \cite{saito2015precision,davis2006relationship}. We use scikit-learn's \texttt{average\_precision\_score} for this calculation. For context, we also report AUC-ROC using the \texttt{roc\_auc\_score} function.

\paragraph{Operational Metric: Precision@K}
Precision@K measures the accuracy of the model within the top-K highest-scored predictions. This is an essential operational metric because it simulates the priority review queue that a fraud analyst would actually use. The system sorts all listings by their predicted fraud probability and then computes what fraction of the top-K listings are true positives. We report Precision@50, @100, and @200, which correspond to the typical daily review capacities of analyst teams.

\paragraph{Other Metrics}
In addition to these core measures, we also track:
\begin{itemize}
    \item \textbf{Precision and Recall}: Standard measures of positive predictive value and sensitivity at a fixed 0.5 threshold.
    \item \textbf{F1-Score}: The harmonic mean of precision and recall.
    \item \textbf{False Positive Rate}: The proportion of legitimate listings that are incorrectly flagged by the model.
\end{itemize}

\subsection{Expanding Window Validation}
\label{sec:4_8_2}

To test how our models handle the evolution of fraud over time, we use an \texttt{ExpandingWindowPipeline}. This approach allows us to systematically validate performance across multiple training periods.

\paragraph{Window Configuration}
The pipeline is configured with a series of increasing windows. For example, it might start by training on 100,000 samples and testing on the following month, then expand to 200,000, 300,000, and finally 400,000 samples.

\paragraph{Pipeline Execution}
The pipeline iterates through these configurations, retrieving the correct temporal slices of data for each step. For each window, it trains the model and then evaluates it on the subsequent period, recording a full set of metrics. This process allows us to identify the minimum amount of data required for stable performance and helps us determine the optimal frequency for retraining the model in a production environment.

\subsection{Drift Detection Algorithms}
\label{sec:4_8_3}

We have implemented two drift detection algorithms to monitor the performance of our models in real-time. While recent research has explored unsupervised and deep learning-based drift detection \cite{arora2024drift}, we chose to use the well-established DDM and ADWIN methods for their proven reliability and ease of interpretation.

\paragraph{DDM (Drift Detection Method)}
DDM monitors the evolution of the model's error rate using a binomial distribution model \cite{gama2004learning}. It triggers a warning or a drift alert whenever the error rate deviates significantly from its historical baseline. The logic for this is outlined in Algorithm~\ref{alg:ddm}.

\begin{algorithm}[H]
\caption{DDM Drift Detection}
\label{alg:ddm}
\begin{algorithmic}[1]
\REQUIRE Stream of predictions $\hat{y}_t$ and true labels $y_t$, warning threshold $\theta_w$, drift threshold $\theta_d$
\ENSURE Drift status $\in \{\text{STABLE}, \text{WARNING}, \text{DRIFT}\}$
\STATE Initialize error history: $\mathcal{E} \gets []$
\FOR{each prediction $(\hat{y}_t, y_t)$ in stream}
    \STATE Compute prediction error: $e_t \gets \mathbb{I}[\hat{y}_t \neq y_t]$
    \STATE Append to history: $\mathcal{E} \gets \mathcal{E} \cup \{e_t\}$
    \STATE Compute error rate: $p \gets \frac{1}{|\mathcal{E}|}\sum_{e \in \mathcal{E}} e$
    \STATE Compute standard deviation: $s \gets \sqrt{\frac{p(1-p)}{|\mathcal{E}|}}$
    \STATE Compute drift statistic: $\delta \gets p + s$
    \IF{$\delta \geq \theta_d$}
        \STATE Reset error history: $\mathcal{E} \gets []$
        \RETURN DRIFT
    \ELSIF{$\delta \geq \theta_w$}
        \RETURN WARNING
    \ELSE
        \RETURN STABLE
    \ENDIF
\ENDFOR
\end{algorithmic}
\end{algorithm}

\paragraph{ADWIN (Adaptive Windowing)}
ADWIN maintains a sliding window of recent errors and detects drift when the distribution of errors between two sub-windows changes significantly \cite{bifet2007learning}. This method is particularly good at adapting the size of the window to the current rate of change in the data.

\begin{algorithm}[H]
\caption{ADWIN Drift Detection}
\label{alg:adwin}
\begin{algorithmic}[1]
\REQUIRE Stream of predictions $\hat{y}_t$ and true labels $y_t$, confidence level $\delta$
\ENSURE Drift status $\in \{\text{STABLE}, \text{DRIFT}\}$
\STATE Initialize sliding window: $\mathcal{W} \gets []$
\FOR{each prediction $(\hat{y}_t, y_t)$ in stream}
    \STATE Compute prediction error: $e_t \gets \mathbb{I}[\hat{y}_t \neq y_t]$
    \STATE Append to window: $\mathcal{W} \gets \mathcal{W} \cup \{e_t\}$
    \FOR{each split point $i = 1$ to $|\mathcal{W}| - 1$}
        \STATE Define sub-windows: $\mathcal{W}_1 \gets \mathcal{W}[1:i]$, $\mathcal{W}_2 \gets \mathcal{W}[i+1:|\mathcal{W}|]$
        \STATE Compute means: $\mu_1 \gets \frac{1}{|\mathcal{W}_1|}\sum_{e \in \mathcal{W}_1} e$, $\mu_2 \gets \frac{1}{|\mathcal{W}_2|}\sum_{e \in \mathcal{W}_2} e$
        \STATE Compute variance estimate: $\sigma^2 \gets \frac{\mu_1(1-\mu_1)}{|\mathcal{W}_1|} + \frac{\mu_2(1-\mu_2)}{|\mathcal{W}_2|}$
        \STATE Compute threshold: $\epsilon_{cut} \gets \sqrt{\frac{2\sigma^2 \ln(2/\delta)}{m}}$ where $m = \frac{1}{\frac{1}{|\mathcal{W}_1|} + \frac{1}{|\mathcal{W}_2|}}$
        \IF{$|\mu_1 - \mu_2| > \epsilon_{cut}$}
            \STATE Drop old sub-window: $\mathcal{W} \gets \mathcal{W}_2$
            \RETURN DRIFT
        \ENDIF
    \ENDFOR
    \RETURN STABLE
\ENDFOR
\end{algorithmic}
\end{algorithm}

Both of these algorithms are designed to run continuously on the production API, with any drift events being logged to trigger automated retraining workflows.

\subsection{Baseline Comparisons}
\label{sec:4_8_4}

To put our results in context, we compare our hybrid models against two sets of baselines:

\paragraph{Model Baselines}
We use simple versions of Logistic Regression and Random Forest to establish a performance floor. These models allow us to quantify exactly how much value is being added by the more complex gradient boosting and graph-based features.

\paragraph{SEON Production Baseline}
We also compare our results against the existing SEON fraud detection system, which uses its own proprietary scores. To ensure a fair comparison, we evaluate both systems on the same test set using the same metrics and perform statistical significance testing using the DeLong test for AUC. All of these comparisons are logged in MLflow to ensure transparency and reproducibility.
